\documentclass[12pt,a4paper]{article}
\usepackage{fontspec}
\usepackage[russian]{babel}
\usepackage{amsmath,amssymb,amsfonts,amsthm}
\usepackage{geometry}
\usepackage{microtype}

\geometry{margin=2cm}
\defaultfontfeatures{Ligatures=TeX}
\setmainfont{Times New Roman}

\begin{document}

\section*{Ответы к зачету}

\begin{enumerate}
    \item \textbf{2.2 Диф-ие и интегрирование вектор-функций.}\\
    \textbf{Определение 2.2.1} (производной вектор ф-ции)\\
    Пусть $\vec{F} : (\alpha, \beta) \to V$ --- функция со значениями в векторном пр-ве $V$ со скалярным произв. \\
    $\vec{a} = \vec{F}'(t_0)$, $t_0 \in (\alpha, \beta)$, если
    \[
    \lim_{t \to t_0} \frac{\vec{F}(t) - \vec{F}(t_0)}{t - t_0} = \vec{a}.
    \]
    \textbf{Лемма 2.2.2}\\
    $\vec{a} = \vec{F}'(t_0) \implies \vec{F}(t) = \vec{F}(t_0) + \vec{a}(t - t_0) + \vec{O}(t)(t - t_0)$, где $\lim_{t \to t_0} \vec{O}(t) = \vec{0}$.\\
    \textbf{Т-ма 2.2.5} (Св-ва произв. и интегралов)\\
    Пусть $\vec{f}, \vec{g}, \vec{h} : [\alpha, \beta] \to V$ --- непрерыв. вр. функции, $\vec{F}, \vec{G}, \vec{H}$ --- их первообразные, $q : [\alpha, \beta] \to \mathbb{R}$, $Q : [\alpha, \beta] \to \mathbb{R}$ --- первообр. $q$. Тогда:
    \begin{enumerate}
        \item[a)] $\vec{F} + \vec{G}$ --- это первообр. ф-ции $\vec{f} + \vec{g}$.
        \item[b)] $c\vec{F}$ первообр. ф-ции $c\vec{f}$, если $c \in \mathbb{R}$.
        \item[c)] $(\vec{F}, \vec{g})$ --- это первообр. функции $(\vec{f}, \vec{g}) + (\vec{F}, \vec{g}')$.
        \item[d)] $q\vec{F}$ --- это первообр. $q\vec{f} + q'\vec{F}$.
        \item[e)] Если $\dim V \le 3$, то $[\vec{F}, \vec{g}]$ --- это первообр. для $[\vec{f}, \vec{g}] + [\vec{F}, \vec{g}']$.
        \item[f)] Если $F^i(t)$ явл-ся первообразной ф-ции $f^i(t)$, то $\vec{F}$ --- первообр. ф-ции $\vec{f}$.
    \end{enumerate}
    Лемма о производной вектор-функции постоянного модуля: \textbf{нет ответа}.

    \item \textbf{2.3 Вектор скорости и касат. к кривой.}\\
    \textbf{Опр 2.3.1} (Параметризованной кривой в пр-ве $X$) с началом в т. $A$ и концом в т. $B$ наз-ся непрер. отображ. $j : [\alpha, \beta] \to X$, $j(\alpha)=A, j(\beta)=B$.\\
    \textbf{Теорема 2.3.2} (касательная прямая): если кривая задана дифференцируемой вектор-функцией $\gamma(t) = (x(t), y(t), z(t))$, то касательная прямая в точке $A = \gamma(t_0)$ имеет вид
    \[
    (X, Y, Z) = (x(t_0), y(t_0), z(t_0)) + \lambda\,(x'(t_0), y'(t_0), z'(t_0)).
    \]
    Нормальная плоскость: \textbf{нет ответа}.

    \item \textbf{Длина кривой и естественная параметризация.}\\
    \textbf{Длина кривой:}
    \[
    L(\gamma) = \int_{\alpha}^{\beta} \|\gamma'(t)\|\,dt.
    \]
    \textbf{Естественная параметризация:}
    \[
    \left\| \frac{d\tilde{\gamma}}{ds} \right\| = 1.
    \]
    Тогда длина дуги между $A$ и $B$ равна $s_1 - s_0$.

    \item \textbf{Кривизна и кручение кривой. Формулы Френе.}\\
    \textbf{Кривизна, нормаль и бинормаль:}\\
    Пусть $\tilde{\gamma}$ --- естественная параметризация кривой. Обозначим
    \[
    \vec{\tau} = \frac{d\tilde{\gamma}}{ds}.
    \]
    Тогда $\vec{\tau}$ называется единичным касательным вектором. \textbf{Кривизной} кривой называется
    \[
    \kappa = \left\| \frac{d\vec{\tau}}{ds} \right\|.
    \]
    Если $\kappa \neq 0$, то
    \[
    \frac{d\vec{\tau}}{ds} = \kappa \vec{n}, \qquad
    \vec{b} = \vec{\tau} \times \vec{n}.
    \]
    Формулы Френе:
    \[
    \frac{d\vec{\tau}}{ds} = \kappa \vec{n}, \qquad
    \frac{d\vec{n}}{ds} = -\kappa \vec{\tau} + \varkappa \vec{b}, \qquad
    \frac{d\vec{b}}{ds} = -\varkappa \vec{n}.
    \]

    \item \textbf{Условия равенства нулю кривизны и кручения.}\\
    \textbf{Утверждение 2.5.3 (об обращении в 0 кривизны и кручения):}\\
    Пусть $\gamma$ --- регулярная непрерывно дифференцируемая кривая.
    \begin{enumerate}
        \item Если $k \equiv 0$, то все точки кривой $\gamma$ лежат на некоторой прямой.
        \item Если $\varkappa \equiv 0$, то кривая лежат в некоторой плоскости.
    \end{enumerate}

    \item \textbf{Определение параметризации, регулярной параметризации и поверхности в евклидовом пространстве.}\\
    \textbf{Параметризованная кривая} --- непрерывное отображение $\gamma: [\alpha, \beta] \to E$.\\
    \textbf{Гладкая поверхность:} множество $S \subset \mathbb{R}^3$ называется \textbf{гладкой поверхностью}, если для любой точки $p \in S$ найдется окрестность $W \subset \mathbb{R}^3$ и гомеоморфизм $r: V \to S \cap W$ (где $V$ --- элементарная область в $\mathbb{R}^2$), такой что:
    \begin{enumerate}
        \item Вектор-функция $r(u, v)$ является гладкой (дифференцируемой).
        \item Векторы частных производных линейно независимы в каждой точке:
        \[
        \frac{\partial r}{\partial u} \times \frac{\partial r}{\partial v} \neq \mathbf{0}.
        \]
        \item Матрица Якоби отображения $r(u,v)$ имеет максимальный ранг:
        \[
        \operatorname{rank}
        \begin{pmatrix}
        \frac{\partial r^1}{\partial u} & \frac{\partial r^2}{\partial u} & \frac{\partial r^3}{\partial u} \\
        \frac{\partial r^1}{\partial v} & \frac{\partial r^2}{\partial v} & \frac{\partial r^3}{\partial v}
        \end{pmatrix} = 2.
        \]
    \end{enumerate}

    \item \textbf{График функции двух переменных как поверхность. Поверхности вращения.}\\
    \textbf{Теорема 2.6.7.}\\
    Пусть $V \subset \mathbb{R}^2$ --- открытое множество, $f: V \to \mathbb{R}$ --- непрерывное отображение. Тогда:
    \[
    \Gamma_f \subset \mathbb{R}^3 \text{ является поверхностью.}
    \]
    При этом отображение $r: V \to \mathbb{R}^3$, задаваемое равенством $r(u,v) = (u, v, f(u,v))$, является параметризацией $\Gamma_f$ в окрестности каждой её точки.\\
    \textbf{Лемма о поверхности вращения:}\\
    Пусть $\gamma(t)$ --- кривая, $\gamma(t) = (f(t), 0, g(t))$. Тогда поверхность вращения $\gamma$ вокруг оси $Z$ задается уравнением:
    \[
    r(u, v) = (f(u)\cos v, f(u)\sin v, g(u)).
    \]

    \item \textbf{Первая и вторая основные формы поверхности.}\\
    \textbf{Определение 2.9.1}\\
    Пусть $S$ --- поверхность, $p \in S$, $T_p S$ --- касательное пространство к поверхности в точке $p$.
    \textbf{1-я основная (или метрическая) форма} поверхности $S$ в точке $p$ --- это отображение:
    \[
    I_p: T_p S \times T_p S \to \mathbb{R}, \quad I_p(\bar a, \bar b) \stackrel{\text{def}}{=} \langle \bar a, \bar b \rangle.
    \]
    \textbf{Утверждение 2.9.2}\\
    Пусть $r: V \to \mathbb{R}^3$ --- регулярная локальная параметризация поверхности $S$, $p = r(u_0, v_0) \in S$. Векторы $\bar{r}_u = \frac{\partial r}{\partial u}(u_0, v_0)$, $\bar{r}_v = \frac{\partial r}{\partial v}(u_0, v_0)$ образуют базис пространства $T_p S$.\\
    Обозначим коэффициенты первой фундаментальной формы:
    \[
    g_{11} = \langle \bar{r}_u, \bar{r}_u \rangle,\quad g_{12} = g_{21} = \langle \bar{r}_u, \bar{r}_v \rangle,\quad g_{22} = \langle \bar{r}_v, \bar{r}_v \rangle.
    \]
    \textbf{Определение 2.9.4 (2-я основная квадратичная форма поверхности)}\\
    Пусть $S$ --- поверхность, $r$ --- регулярная параметризация, $r(u_0, v_0) = p \in S$. Форма $II_p: T_p S \times T_p S \to \mathbb{R}$ называется \textbf{второй основной билинейной формой}, если она задается равенством:
    \[
    II_p(\bar{a}, \bar{b}) = \sum_{i,j=1}^2 b_{ij} a^i b^j.
    \]
    \textbf{Лемма 2.9.5}\\
    Вторая основная форма поверхности $S$ в точке $p$ при замене параметров либо не меняется, либо умножается на $-1$, т.е. определена с точностью до знака. При этом знак меняется только тогда, когда вектор нормали $\bar{m}$ умножается на $(-1)$.

    \item \textbf{Нормальная кривизна. Теорема Менье.}\\
    \textbf{Теорема 2.9.6 --- Теорема Менье}\\
    Пусть $\gamma(t) \in S$ --- кривая на поверхности $S$, проходящая через точку $p$. Пусть $n$ --- вектор главной нормали кривой в точке $p$, а $m$ --- единичный вектор нормали к самой поверхности $S$ в точке $p$. Тогда выполнено соотношение для нормальной кривизны $k_n$:
    \[
    k_n = (k \cdot n, m) = k \cos \theta = \frac{II_p(\gamma'(t_0))}{I_p(\gamma'(t_0))}.
    \]

    \item \textbf{Главные кривизны и главные направления. Теорема Эйлера.}\\
    \textbf{Теорема 2.9.9 --- Теорема Эйлера}\\
    Пусть $S$ --- регулярная поверхность, $p \in S$. Тогда в касательной плоскости существуют два взаимно перпендикулярных направления $l_1$ и $l_2$, в которых нормальная кривизна принимает экстремальные значения $k_1$ и $k_2$ (главные кривизны).
    Для любого направления $l$, составляющего угол $\theta$ с главным направлением $l_1$, нормальная кривизна $k_l$ вычисляется по формуле:
    \[
    k_l = k_1 \cos^2 \theta + k_2 \sin^2 \theta.
    \]

    \item \textbf{Уравнение для главных кривизн. Гауссова кривизна.}\\
    \textbf{Теорема Гаусса / Theorema Egregium}\\
    Гауссова кривизна выражается только через компоненты первой фундаментальной формы $g_{ij}$ и их частные производные:
    \[
    K_g = \frac{b_{11}b_{22} - b_{12}^2}{g_{11}g_{22} - g_{12}^2} = \frac{\det(II)}{\det(I)}.
    \]
    Уравнение для главных кривизн: \textbf{нет ответа}.

    \item \textbf{Деривационные уравнения. Символы Кристоффеля.}\\
    \textbf{Деривационные уравнения (Гаусса-Вейнгартена)} запишем в матричной форме:
    \[
    \frac{\partial}{\partial u}
    \begin{pmatrix}
    \bar{r}_u \\
    \bar{r}_v \\
    \bar{m}
    \end{pmatrix}
    =
    \begin{pmatrix}
    \Gamma_{11}^1 & \Gamma_{11}^2 & b_{11}\\
    \Gamma_{12}^1 & \Gamma_{12}^2 & b_{12}\\
    -b_1^1 & -b_1^2 & 0
    \end{pmatrix}
    \begin{pmatrix}
    \bar{r}_u \\
    \bar{r}_v \\
    \bar{m}
    \end{pmatrix},
    \]
    \[
    \frac{\partial}{\partial v}
    \begin{pmatrix}
    \bar{r}_u \\
    \bar{r}_v \\
    \bar{m}
    \end{pmatrix}
    =
    \begin{pmatrix}
    \Gamma_{12}^1 & \Gamma_{12}^2 & b_{12}\\
    \Gamma_{22}^1 & \Gamma_{22}^2 & b_{22}\\
    -b_2^1 & -b_2^2 & 0
    \end{pmatrix}
    \begin{pmatrix}
    \bar{r}_u \\
    \bar{r}_v \\
    \bar{m}
    \end{pmatrix}.
    \]
    Символы Кристоффеля: \textbf{нет ответа}.

    \item \textbf{Уравнения Гаусса-Кодацци.}\\
    \textbf{Теорема 2.10.3 (Об уравнениях Гаусса-Кодацци)}\\
    Имеют место следующие уравнения консистентности:
    \[
    \frac{\partial A_2}{\partial u} - \frac{\partial A_1}{\partial v} + [A_1, A_2] = 0.
    \]

    \textbf{Теорема 2.10.4 (Теорема Гаусса / Theorema Egregium)}\\
    Гауссова кривизна выражается только через компоненты первой фундаментальной формы $g_{ij}$ и их частные производные:
    \[
    K_g = \frac{b_{11}b_{22} - b_{12}^2}{g_{11}g_{22} - g_{12}^2} = \frac{\det(II)}{\det(I)}.
    \]

    \item \textbf{Ковариантное дифференцирование. Уравнения геодезических. Геодезические на поверхностях вращения.}\\
    нет ответа

    \item \textbf{Геодезические, как экстремали функционалов действия и функционалов длины.}\\
    нет ответа
\end{enumerate}

\end{document}
